\chapter{Conclusion}

\section{Summary of Work}

This project developed and evaluated a reinforcement learning-based framework for adaptive pilot and power allocation in an OFDM-based Integrated Sensing and Communication (ISAC) system. The work addressed the three-way tradeoff between communication spectral efficiency, synchronization accuracy, and sensing precision that arises from the dual role of pilot subcarriers in OFDM-ISAC.

The key contributions delivered are:

\begin{enumerate}
    \item \textbf{OFDM-ISAC Simulation Environment:} A frame-level discrete-event simulator was implemented incorporating mean-reverting AR(1) CFO dynamics, correlated block-fading channel evolution via an AR(1) process, and AWGN. The simulator computes physically consistent receiver-observable metrics at each frame without knowledge of the true channel state.

    \item \textbf{Observable Estimation Metrics:} A least-squares phase-difference CFO estimator and an observable synchronization mismatch metric $\cE_t$ were derived and implemented. These quantities are computable directly from received pilot symbols and serve as both state variables and reward signal components.

    \item \textbf{Empirical Sensing Quality Metric:} The sensing metric $S_t$ was defined as the sidelobe-to-peak ratio of the pilot-domain delay-correlation profile. This metric is fully computable at the receiver without access to ground-truth channel parameters and provides a meaningful surrogate for sensing resolution.

    \item \textbf{MDP Formulation:} The joint pilot ratio and pilot power fraction selection problem was formulated as a finite-horizon MDP with a seven-dimensional observable state, twelve discrete actions, and a normalized three-component reward function.

    \item \textbf{DQN Agent:} A Deep Q-Network with a two-hidden-layer MLP architecture (48 neurons per layer), experience replay, and a target network was trained to select $(\rho_t, \beta_t)$ at each frame. The implementation is self-contained in NumPy and requires no deep-learning framework dependency.

    \item \textbf{Structured Reward-Weight Selection:} A stability-aware multi-seed sweep methodology was used to select the reward weights $(\lambda_{\text{rate}}, \lambda_{\text{sync}}, \lambda_{\text{sensing}}) = (1.0, 0.2, 0.4)$ from three confirmed candidates, providing a reproducible and defensible configuration.

    \item \textbf{Baseline Comparison:} The trained DQN was compared against twelve fixed policies, a random policy, and a heuristic policy. The DQN outperformed the random policy (by $8.20$ reward units) and the heuristic policy (by $14.34$ reward units), confirming that a nontrivial adaptive policy was learned.
\end{enumerate}

\section{Limitations}

The results should be interpreted in the context of several limitations:

\begin{itemize}
    \item \textbf{Stationary environment:} The channel and CFO dynamics are governed by fixed AR(1) parameters throughout training and evaluation. Under a stationary distribution, a well-chosen constant allocation can match or exceed an adaptive policy. The adaptive policy's advantage only materialises in non-stationary environments where the optimal action changes over time.

    \item \textbf{Hindsight-best fixed baseline:} The best fixed policy ($\rho{=}0.1, \beta{=}0.2$, reward $53.05$) was identified after observing experimental results. In practice, the optimal fixed allocation is not known in advance and must itself be tuned. The comparison is therefore intentionally strict.

    \item \textbf{Short training:} The DQN was trained for only 100 episodes of 64 frames each ($6400$ environment steps), which is modest for a tabular-equivalent problem of this size. An extended 500-episode run with the same reward weights was also carried out and showed a clear reward plateau near $45$ after episode 100, indicating that the ceiling observed in the 100-episode protocol is not purely a consequence of under-training; further progress would require a richer environment rather than simply longer training.

    \item \textbf{Flat-fading assumption:} The system model assumes the channel is flat across all $K$ subcarriers within a frame. This simplification is reasonable for narrowband ISAC but does not capture frequency-selective multipath effects relevant to wideband systems.

    \item \textbf{SISO operation:} The framework is limited to a single transmit and receive antenna. Extension to MIMO would require joint beamforming and pilot allocation design.

    \item \textbf{Python-level simulator loop:} While the DQN agent itself is vectorised in NumPy, the per-frame OFDM-ISAC simulator still executes through a Python-level loop over subcarriers and symbols. For the present problem size this is acceptable, but scaling to wider bandwidths or frequency-selective channels would benefit from vectorising or JIT-compiling the simulator's inner loop.
\end{itemize}

\section{Future Work}

The following extensions would meaningfully advance the system beyond the current proof-of-concept:

\begin{enumerate}
    \item \textbf{Non-stationary environment with regime switching:} Modelling the CFO process as switching between a low-drift and high-drift regime would create a genuinely dynamic problem where adaptive control is necessary. This is the most direct path to demonstrating clear DQN superiority over fixed baselines.

    \item \textbf{Vectorised simulator inner loop:} Refactoring the simulator's per-subcarrier and per-symbol loops into NumPy array operations (or compiling them via Numba / JAX) would further reduce the cost per training step and enable wider hyperparameter sweeps within practical time budgets.

    \item \textbf{Constrained reward formulation:} Rather than combining objectives into a scalar reward, a constrained MDP formulation could enforce hard minimum thresholds on synchronization quality while maximizing rate. This would provide cleaner separation of communication and sensing guarantees.

    \item \textbf{Frequency-selective channel model:} Extending the channel from flat-fading to a tapped-delay-line multipath model would make the delay-domain sensing metric more physically meaningful and introduce subcarrier-dependent channel estimation requirements.

    \item \textbf{Hardware prototyping:} Implementing the system on a software-defined radio (SDR) platform would validate the observable metrics and reward function under real hardware impairments including phase noise, nonlinear amplification, and timing jitter.

    \item \textbf{Multi-user extension:} Extending the problem to multiple users sharing the same spectrum would introduce inter-user interference as an additional dimension of the pilot allocation tradeoff.
\end{enumerate}

\section{Closing Remarks}

This project demonstrates that a DQN-based adaptive pilot allocation policy can be trained from receiver-observable metrics in an OFDM-ISAC simulator. The framework is physically grounded, causally consistent, and honestly evaluated. The honest identification of the current performance ceiling and its root cause — the stationary environment and the dominance of the rate term — is itself a contribution, as it clearly motivates the next steps needed to transform this proof-of-concept into a competitive adaptive system.
